\section{Implémentation et Résultats}
\label{sec:task3}

\subsection{Comparatif des Diodes}
Pour maximiser le rendement RF-DC à très faible puissance, nous avons comparé différentes diodes Schottky. Le package \texttt{booktabs} permet d'illustrer ces résultats clairement (Tableau \ref{tab:diodes}).

\begin{table}[h]
    \centering
    \caption{Caractéristiques des diodes Schottky testées}
    \label{tab:diodes}
    \begin{tabular}{lcccc}
        \toprule
        \textbf{Diode} & \textbf{($I_s$)} & \textbf{($R_s$)} & \textbf{($C_j$)} & \textbf{($F$)}\\
        \midrule
        HSMS-2820 & \SI{22}{\nano\ampere} & \SI{6}{\ohm} & \SI{0.7}{\pico\farad} & up to \SI{4}{\giga\hertz} \\
         HSMS-2850 & \SI{3000}{\nano\ampere} & \SI{25}{\ohm} & \SI{0.18}{\pico\farad} & below \SI{1.5}{\giga\hertz} \\
        HSMS-2860   & \SI{50}{\nano\ampere} & \SI{6}{\ohm} & \SI{0.18}{\pico\farad} & up to \SI{5.8}{\giga\hertz} \\
        \bottomrule
    \end{tabular}
\end{table}

\subsection{Script de calcul de bilan de liaison}
Afin d'automatiser nos calculs de densité de puissance, nous avons développé un petit script Python.

\begin{lstlisting}[language=Python, caption=Conversion dBm en Watts]
def dbm_to_watts(dbm_value):
    # Convertit une puissance reçue en dBm vers des Watts
    milliwatts = 10 ** (dbm_value / 10)
    watts = milliwatts / 1000
    return watts

puissance_rx = -20 # dBm
print(f"Puissance : {dbm_to_watts(puissance_rx)} W")
\end{lstlisting}